\paragraph{潜在生産水準目標（POLT）}
\label{sec:chapter5_a_shock_policies_polt}
POLTの式については
\ref{sec:chapter5_beta_shock_policies_polt}
を参照のこと。
\begin{enumerate}
    \item
    \label{itm:chapter5_a_shock_policies_polt_c_H_W}
        \(c_t^{H\to W}\)：
        総消費指数\(c_t^{H\to W}\)の図\ref{fig:chapter5_a_shock_graphs_c_H_W}を見ると、POLTの軌道は、供給ショックに伴う潜在的な生産能力の低下を反映し、ショック直後の初期においてマイナス方向への下落を引き起こしていることがわかる。
        生産者物価水準目標（PPLT）や消費者物価インフレ目標（IT-CPI）などで見られた実体経済への過度な下押し圧力による極端な大暴落とは異なり、POLTは目標とする生産水準そのものがショックに応じて低下するため、経済の基礎的条件に沿った形での消費の減少となっている。
        しかしながら、消費の絶対的な低下自体は厚生にマイナスに働くため、この前半における消費の落ち込みがPOLTの全体的な厚生評価を押し下げる要因となっている。
        また深く落ち込んだ状態から、ショックの減衰に伴う潜在水準の回復に合わせて長い時間をかけて定常状態へ回帰する軌道を描いている。
    \item
    \label{itm:chapter5_a_shock_policies_polt_y_H}
        \(y_t^H\)：
        生産\(y_t^H\)の図\ref{fig:chapter5_a_shock_graphs_y_H}を見ると、POLTの軌道は総消費指数\(c_t^{H\to W}\)と同様に、ショック直後の初期においてマイナス方向へ落ち込んでいることがわかる。
        しかしながら、生産者物価水準目標（PPLT）や消費者物価インフレ目標（IT-CPI）などで見られたような極端な大暴落とは軌道が異なっている。POLTは供給ショックによって低下した潜在生産水準そのものを目標として追随するため、生産の落ち込みは経済の基礎的条件の変化（生産性の低下）に沿ったものに留められている。
        生産の低下は労働の減少を意味するため、この前半部分においては労働の非効用を減らして厚生にプラスに働く効果をもたらしている。
        その後は、ショックの減衰に伴って潜在生産水準が回復していくのに合わせ、長い時間をかけて定常状態へ回帰する軌道を描いている。
    \item
    \label{itm:chapter5_a_shock_policies_polt_pi_H}
        \(\pi_t^H\)：
        グロス・インフレ率\(\pi_t^H\)の図\ref{fig:chapter5_a_shock_graphs_pi_H}を見ると、POLTの軌道は生産者物価水準目標（PPLT）や生産者物価インフレ目標（IT-PPI）のように全期間を通じてゼロ近傍で完全に平坦な推移を維持するわけではなく、ショックに応じてある程度の変動を許容していることがわかる。
        これは、POLTが物価水準やインフレ率そのものの安定化を直接の目標とするのではなく、実体経済（生産）を変化した潜在生産水準に追随させることを最優先とする政策であるためである。
        消費者物価インフレ目標（IT-CPI）などで見られたような極端なマイナス方向（デフレ）への下落とは異なり、POLTにおけるインフレ率の変動は、供給ショックに伴う相対価格の調整を促し、生産を目標水準へと導くための必要なプロセスとして機能している。しかしながら、このインフレ率の変動が次項で述べる価格分散の発生を招く要因にもなっている。
    \item
    \label{itm:chapter5_a_shock_policies_polt_Delta_t_minus_1_H}
        \(\Delta_{t-1}^H\)：
        価格分散\(\Delta_{t-1}^H\)の図\ref{fig:chapter5_a_shock_graphs_Delta_H}を見ると、POLTの定常状態からの乖離は、生産者物価水準目標（PPLT）や生産者物価インフレ目標（IT-PPI）が全期間を通じてほぼゼロの水準を完全に維持していることと比較して、明確な増大を見せていることがわかる。
        これは前述のインフレ率の変動を許容した結果として、価格改定のタイミングの違いに伴う相対価格のばらつきが生じたためである。
        この価格分散の蓄積は生産効率の悪化を意味しており、POLTの厚生を直接的に押し下げるマイナス要因として働いている。
        POLTは実体経済の不必要な大暴落を回避し、生産を潜在水準に沿って調整することには成功しているものの、その代償として生じるこの価格分散による効率性の喪失が、厚生上の大きな制約となっていると言える。
\end{enumerate}